\usepackage{graphicx} % 用于插入图片
\usepackage{setspace} % 用于设置行距
\usepackage{color}
\usepackage{xcolor}
\usepackage[absolute, overlay]{textpos}
\usepackage{tcolorbox} % 用于创建带边框的盒子
\usepackage{fancyhdr}
\tcbuselibrary{skins} % 加载 skins 库以获得更多样式选项

\usepackage[
  style=gb7714-2015ay,
  backend=biber,
  maxcitenames=2,                 % cite 时最多列 2 名作者；超过即用 et al
  mincitenames=1,                 % et al 时只保留第 1 名作者
  maxbibnames=20,                 % 参考文献列表最多列 20 名作者
  minbibnames=20,
  gbnamefmt=lowercase,            % 提交稿：Goodfellow I，而非全大写姓氏
  gbfieldtype=false,             % 提交稿常见著录不带[M]/[J]标识
  gbpub=false,
  sortcites=false,                % cite 内部不自动排序（在源码里按年代写）
  sorting=gb7714-2015             % 参考文献列表：中文在前外文在后、按姓氏
]{biblatex}
\DeclareDelimFormat[cite,parencite,pagescite,citep]{nameyeardelim}{\addspace}
\DeclareDelimFormat[cite,parencite,textcite]{finalnamedelim}{%
  \iffieldequalstr{langid}{chinese}{和}{\addspace and\addspace}}
% 规范§7.5+§8.1要求：et al 后不加英文点（参考文献列表和正文引用都一致）
% gb7714-2015ay 风格 cite 内部用 andothersincite，列表内部用 andothers
\DefineBibliographyStrings{english}{%
  andothers       = {et al},%
  andothersincite = {et al},%
}
% 规范要求 2 名作者用 "Smith and Jones 1992" 形式（用 and，不用逗号）
% 注意：gb7714-2015ay 风格里参考文献列表用顿号/逗号是规范的，
% 但 cite 时按"姓名加年份"规则需要 and。
\renewcommand*{\multinamedelim}{\addcomma\space}
\renewcommand*{\finalnamedelim}{\addspace and\addspace}
\AtBeginBibliography{%
  \renewcommand*{\multinamedelim}{\addcomma\space}%
  \renewcommand*{\finalnamedelim}{\addcomma\space}%
}

% 多篇文献之间用逗号，而不是分号
\DeclareDelimFormat{multicitedelim}{\addcomma\space}
\DeclareDelimFormat{compcitedelim}{\addcomma\space}
\DeclareFieldFormat[article,inproceedings,incollection,unpublished]{title}{#1}
\addbibresource{\thesisBibFile}  % 替换为你的 .bib 文件名
% Automatic common entry formats modelled on both submitted thesis reference lists.
% Less common types keep the package's fallback format and need individual review.
\DeclareFieldFormat{title}{#1}
\DeclareFieldFormat[book,article,inproceedings,incollection,inbook,thesis,report,online]{title}{#1}
\DeclareFieldFormat{booktitle}{#1}
\DeclareFieldFormat{journaltitle}{\mkbibemph{#1}}
\DeclareFieldFormat{pages}{#1}
\DeclareFieldFormat{url}{\url{#1}}
\DeclareFieldFormat{doi}{\url{https://doi.org/#1}}
\newcommand*{\hzauBibPeriod}{\iffieldequalstr{langid}{chinese}{\unspace．}{\addperiod\space}}
\newcommand*{\hzauBibComma}{\iffieldequalstr{langid}{chinese}{\unspace，}{\addcomma\space}}
\newcommand*{\hzauBibColon}{\iffieldequalstr{langid}{chinese}{\unspace：}{\addcolon\space}}
\newbibmacro*{hzau:year}{\printfield{year}\printfield{extradate}}
\newbibmacro*{hzau:author-title}{%
  \usebibmacro{bibindex}\usebibmacro{begentry}%
  \ifnameundef{author}{\printnames{editor}}{\printnames{author}}%
  \setunit{\hzauBibPeriod}\printfield{title}}
\newbibmacro*{hzau:links}{%
  \iffieldundef{doi}{%
    \iffieldundef{url}{}{\setunit{\addperiod\space}\printfield{url}}%
  }{\setunit{\addperiod\space}\printfield{doi}}}
\AtBeginBibliography{\renewcommand*{\finentrypunct}{}}

\DeclareBibliographyDriver{book}{%
  \usebibmacro{hzau:author-title}%
  \iffieldundef{edition}{}{\setunit{\hzauBibPeriod}\printfield{edition}}%
  \setunit{\hzauBibPeriod}\printlist{location}%
  \setunit*{\hzauBibColon}\printlist{publisher}%
  \setunit{\hzauBibComma}\usebibmacro{hzau:year}%
  \iffieldundef{pages}{}{\setunit{\hzauBibPeriod}\printfield{pages}}%
  \usebibmacro{hzau:links}\usebibmacro{finentry}}

\DeclareBibliographyDriver{article}{%
  \usebibmacro{hzau:author-title}%
  \setunit{\hzauBibPeriod}\printfield{journaltitle}%
  \setunit{\hzauBibComma}\usebibmacro{hzau:year}%
  \iffieldundef{volume}{}{\setunit{\hzauBibComma}\printfield{volume}}%
  \iffieldundef{number}{}{\printtext{(\printfield{number})}}%
  \iffieldundef{pages}{}{\setunit{\hzauBibColon}\printfield{pages}}%
  \usebibmacro{hzau:links}\usebibmacro{finentry}}

\DeclareBibliographyAlias{inproceedings}{hzau-proceedings}
\DeclareBibliographyDriver{hzau-proceedings}{%
  \usebibmacro{hzau:author-title}%
  \setunit{\hzauBibPeriod}\printfield{booktitle}%
  \setunit{\hzauBibComma}\usebibmacro{hzau:year}%
  \iflistundef{location}{}{\setunit{\hzauBibComma}\printlist{location}}%
  \iffieldundef{pages}{}{\setunit{\hzauBibComma}\printfield{pages}}%
  \usebibmacro{hzau:links}\usebibmacro{finentry}}

\DeclareBibliographyAlias{incollection}{hzau-collection}
\DeclareBibliographyDriver{hzau-collection}{%
  \usebibmacro{hzau:author-title}%
  \setunit{\hzauBibPeriod}\printtext{\iffieldequalstr{langid}{chinese}{见：}{In:\space}}%
  \ifnameundef{editor}{}{\printnames{editor}\setunit{\addcomma\space}}%
  \printfield{booktitle}%
  \setunit{\hzauBibPeriod}\printlist{location}%
  \setunit*{\hzauBibColon}\printlist{publisher}%
  \setunit{\hzauBibComma}\usebibmacro{hzau:year}%
  \iffieldundef{pages}{}{\setunit{\hzauBibPeriod}\printfield{pages}}%
  \usebibmacro{hzau:links}\usebibmacro{finentry}}

\DeclareBibliographyAlias{thesis}{hzau-thesis}
\DeclareBibliographyDriver{hzau-thesis}{%
  \usebibmacro{hzau:author-title}%
  \setunit{\hzauBibPeriod}\printfield{type}%
  \setunit{\hzauBibPeriod}\printlist{location}%
  \setunit*{\hzauBibColon}\printlist{institution}%
  \setunit{\hzauBibComma}\usebibmacro{hzau:year}%
  \usebibmacro{hzau:links}\usebibmacro{finentry}}

\DeclareBibliographyDriver{online}{%
  \usebibmacro{hzau:author-title}%
  \setunit{\hzauBibPeriod}\usebibmacro{hzau:year}%
  \usebibmacro{hzau:links}\usebibmacro{finentry}}


% 给参考文献列表加顺序编号（规范例子里"1.姜锡州..." "2.蓝盛银..."这种"序号."风格）
% 因为 gb7714-2015ay 风格不设置 labelnumber 字段，所以用计数器方式自增
\newcounter{bibrefcnt}
\defbibenvironment{bibliography}
  {\setcounter{bibrefcnt}{0}%
   \list{\stepcounter{bibrefcnt}\arabic{bibrefcnt}.}%
        {\setlength{\labelwidth}{2.2em}%
         \setlength{\leftmargin}{\labelwidth}%
         \setlength{\labelsep}{0.5em}%
         \addtolength{\leftmargin}{\labelsep}%
         \setlength{\itemsep}{0.4ex plus 0.2ex}%
         \setlength{\parsep}{0pt}%
         \renewcommand*{\makelabel}[1]{##1\hss}}}%
  {\endlist}%
  {\item}

% 自定义参考文献标题格式（居中+目录）
\defbibheading{bibliography}[参考文献]{%
  \clearpage\phantomsection
  \addcontentsline{toc}{section}{参考文献}%
  \vspace*{-10pt}%
  \noindent\makebox[\textwidth][c]{\heiti\zihao{3}参\hspace{1em}考\hspace{1em}文\hspace{1em}献}\par%
  \vspace{8pt}%
}


\definecolor{lightgray}{rgb}{.9,.9,.9}
\definecolor{darkgray}{rgb}{.4,.4,.4}
\definecolor{purple}{rgb}{0.65, 0.12, 0.82}
\definecolor{hzaugreen}{HTML}{67b356} 

%%%%%%%%%%%%%%%%
\usepackage{bicaption} % 支持双标题
% 设置中英文标题分隔符（默认为换行）
\DeclareCaptionOption{separator}{\def\captionseparator{#1}}
\captionsetup[bi-first]{name=图} % 中文"图"代替"Figure"
\captionsetup[bi-second]{name=Figure}
% 设置表格中英文标题格式
\captionsetup[table][bi-first]{name=表} % 中文"表"代替"Table"
\captionsetup[table][bi-second]{name=Table}
% 设置 longtable 中英文标题格式
\captionsetup[longtable][bi-first]{name=表}
\captionsetup[longtable][bi-second]{name=Table}



% The bibliography path is set by the entry file before this file is loaded.
% Replace the earlier resource declaration with the selected path.
